% GENERATED FILE. Edit canonical lessons, then run npm run generate:books.
% canonical-source-hash: fnv1a64:010e97520f3c1171
% canonical-lessons: RU-C476-ispug, RU-C476-pugatsya, RU-C476-vzdragivat, RU-C476-trus, RU-C476-smelost, RU-R476-review-469-476

\chapter{Fright and Courage}
\label{ch:ru-fright-and-courage}
\hlchaptermodality{476}

\begin{chapteropening}
\textbf{By the end of this chapter:} \emph{I can talk about getting a fright and being brave.}

Complete the last lesson of chapter 476: I can talk about getting a fright and being brave.
\end{chapteropening}

\section[\texorpdfstring{\ru{испуг} (ispúg)}{испуг (ispúg)}]{\ru{испуг} (ispúg) — a fright}

\label{lesson:RU-C476-ispug}



\begin{quote}
Before the new one: say the Russian for pride, then the Russian for to boast.
\end{quote}

\subsection*{You'll want to know: \ru{испуг}}
\textbf{\ru{испуг}} (\emph{ispúg}) — \textquotedblleft{}a fright\textquotedblright{}.

\begin{quote}
\ru{От} \ru{испуга} \ru{кошка} \ru{прыгнула} \ru{на} \ru{шкаф}.
\end{quote}

\begin{quote}
\emph{Ot ispúga kóshka prýgnula na shkaf.}
\end{quote}

\begin{quote}
The cat jumped onto the wardrobe in fright.
\end{quote}

Masculine. The adjective is \emph{\ru{испуганный}} — frightened.

\begin{cousinweb}
\emph{\ru{испуг}} is built on \emph{\ru{ис}-} and the root of \emph{\ru{пугать}} 'to frighten'; EtymDB lists \emph{\ru{пугать}} without a source, so its further origin is not settled.
\end{cousinweb}

\begin{grammarlens}[title={What you've built}]
One more word for telling, explaining and describing at B1.
\end{grammarlens}

\subsection*{Your turn}
\begin{itemize}
\raggedright
  \item \emph{Say it:} \emph{ispúg}
  \item \emph{Say it:} \emph{ispúg}, once more
  \item \emph{Say it:} a sentence of your own with \emph{ispúg}
  \item \emph{Recall:} say the Russian for pride, then the Russian for to boast, then say \emph{ispúg} again
\end{itemize}

\begin{tcolorbox}[breakable,colback=teal!4,colframe=teal!35!black,title={Before you move on}]
What does \ru{испуг} mean? (\textquotedblleft{}A fright\textquotedblright{}.) Say it once more.
\end{tcolorbox}
\section[\texorpdfstring{\ru{пугаться} (pugátsya)}{пугаться (pugátsya)}]{\ru{пугаться} (pugátsya) — to get frightened}

\label{lesson:RU-C476-pugatsya}



\begin{quote}
Before the new one: say the Russian for to boast, then the Russian for a fright.
\end{quote}

\subsection*{You'll want to know: \ru{пугаться}}
\textbf{\ru{пугаться}} (\emph{pugátsya}) — \textquotedblleft{}to get frightened\textquotedblright{}.

\begin{quote}
\ru{Мой} \ru{сын} \ru{пугается} \ru{грома}.
\end{quote}

\begin{quote}
\emph{Moy syn pugáyetsya gróma.}
\end{quote}

\begin{quote}
My son gets frightened of thunder.
\end{quote}

Imperfective; perfective \ru{испугаться}. Takes the genitive: \emph{\ru{пугаться} \ru{грома}}.

\begin{cousinweb}
\emph{\ru{пугаться}} is the reflexive of \emph{\ru{пугать}} 'to frighten, to scare'; EtymDB lists \emph{\ru{пугать}} without a source, so the origin is not settled.
\end{cousinweb}

\begin{grammarlens}[title={What you've built}]
One more word for telling, explaining and describing at B1.
\end{grammarlens}

\subsection*{Your turn}
\begin{itemize}
\raggedright
  \item \emph{Say it:} \emph{pugátsya}
  \item \emph{Say it:} \emph{pugátsya}, once more
  \item \emph{Say it:} a sentence of your own with \emph{pugátsya}
  \item \emph{Recall:} say the Russian for to boast, then the Russian for a fright, then say \emph{pugátsya} again
\end{itemize}

\begin{tcolorbox}[breakable,colback=teal!4,colframe=teal!35!black,title={Before you move on}]
What does \ru{пугаться} mean? (\textquotedblleft{}To get frightened\textquotedblright{}.) Say it once more.
\end{tcolorbox}
\section[\texorpdfstring{\ru{вздрагивать} (vzdrágivat)}{вздрагивать (vzdrágivat)}]{\ru{вздрагивать} (vzdrágivat) — to flinch}

\label{lesson:RU-C476-vzdragivat}



\begin{quote}
Before the new one: say the Russian for a fright, then the Russian for to get frightened.
\end{quote}

\subsection*{You'll want to know: \ru{вздрагивать}}
\textbf{\ru{вздрагивать}} (\emph{vzdrágivat}) — \textquotedblleft{}to flinch\textquotedblright{}.

\begin{quote}
\ru{Собака} \ru{вздрагивает} \ru{от} \ru{каждого} \ru{звука}.
\end{quote}

\begin{quote}
\emph{Sobáka vzdrágivayet ot kázhdogo zvúka.}
\end{quote}

\begin{quote}
The dog flinches at every sound.
\end{quote}

Imperfective; perfective \ru{вздрогнуть} (once). Related to \emph{\ru{дрожать}} — to tremble.

\begin{cousinweb}
\emph{\ru{вздрагивать}} is built on \emph{\ru{вз}-} and the root of \emph{\ru{дрогнуть}} 'to quiver' and \emph{\ru{дрожать}} 'to tremble'; EtymDB lists these words without a source.
\end{cousinweb}

\begin{grammarlens}[title={What you've built}]
One more word for telling, explaining and describing at B1.
\end{grammarlens}

\subsection*{Your turn}
\begin{itemize}
\raggedright
  \item \emph{Say it:} \emph{vzdrágivat}
  \item \emph{Say it:} \emph{vzdrágivat}, once more
  \item \emph{Say it:} a sentence of your own with \emph{vzdrágivat}
  \item \emph{Recall:} say the Russian for a fright, then the Russian for to get frightened, then say \emph{vzdrágivat} again
\end{itemize}

\begin{tcolorbox}[breakable,colback=teal!4,colframe=teal!35!black,title={Before you move on}]
What does \ru{вздрагивать} mean? (\textquotedblleft{}To flinch\textquotedblright{}.) Say it once more.
\end{tcolorbox}
\section[\texorpdfstring{\ru{трус} (trus)}{трус (trus)}]{\ru{трус} (trus) — a coward}

\label{lesson:RU-C476-trus}



\begin{quote}
Before the new one: say the Russian for to get frightened, then the Russian for to flinch.
\end{quote}

\subsection*{You'll want to know: \ru{трус}}
\textbf{\ru{трус}} (\emph{trus}) — \textquotedblleft{}a coward\textquotedblright{}.

\begin{quote}
\ru{В} \ru{начале} \ru{фильма} \ru{герой} \ru{был} \ru{трусом}.
\end{quote}

\begin{quote}
\emph{V nachále fílma geróy byl trúsom.}
\end{quote}

\begin{quote}
At the start of the film the hero was a coward.
\end{quote}

Masculine. The adjective is \emph{\ru{трусливый}} — cowardly.

\begin{cousinweb}
EtymDB lists \emph{\ru{трус}} without a source; it is usually linked to \emph{\ru{трусить}} 'to shake, to tremble', so a coward is someone who trembles.
\end{cousinweb}

\begin{grammarlens}[title={What you've built}]
One more word for telling, explaining and describing at B1.
\end{grammarlens}

\subsection*{Your turn}
\begin{itemize}
\raggedright
  \item \emph{Say it:} \emph{trus}
  \item \emph{Say it:} \emph{trus}, once more
  \item \emph{Say it:} a sentence of your own with \emph{trus}
  \item \emph{Recall:} say the Russian for to get frightened, then the Russian for to flinch, then say \emph{trus} again
\end{itemize}

\begin{tcolorbox}[breakable,colback=teal!4,colframe=teal!35!black,title={Before you move on}]
What does \ru{трус} mean? (\textquotedblleft{}A coward\textquotedblright{}.) Say it once more.
\end{tcolorbox}
\section[\texorpdfstring{\ru{смелость} (smélost)}{смелость (smélost)}]{\ru{смелость} (smélost) — courage, boldness}

\label{lesson:RU-C476-smelost}



\begin{quote}
Before the new one: say the Russian for to flinch, then the Russian for a coward.
\end{quote}

\subsection*{You'll want to know: \ru{смелость}}
\textbf{\ru{смелость}} (\emph{smélost}) — \textquotedblleft{}courage, boldness\textquotedblright{}.

\begin{quote}
\ru{Нужно} \ru{много} \ru{смелости}, \ru{чтобы} \ru{петь} \ru{на} \ru{сцене}.
\end{quote}

\begin{quote}
\emph{Núzhno mnógo smélosti, chtóby pet na stséne.}
\end{quote}

\begin{quote}
It takes a lot of courage to sing on stage.
\end{quote}

Feminine. Built on \emph{\ru{смелый}} — bold, brave.

\begin{cousinweb}
\emph{\ru{смелость}} is built on \emph{\ru{смелый}} 'bold' with \emph{-\ru{ость}} (EtymDB); \emph{\ru{смелый}} is inherited from Old East Slavic \emph{s\ru{ъ}měl\ru{ъ}} (EtymDB).
\end{cousinweb}

\begin{grammarlens}[title={What you've built}]
One more word for telling, explaining and describing at B1.
\end{grammarlens}

\subsection*{Your turn}
\begin{itemize}
\raggedright
  \item \emph{Say it:} \emph{smélost}
  \item \emph{Say it:} \emph{smélost}, once more
  \item \emph{Say it:} a sentence of your own with \emph{smélost}
  \item \emph{Recall:} say the Russian for to flinch, then the Russian for a coward, then say \emph{smélost} again
\end{itemize}

\begin{tcolorbox}[breakable,colback=teal!4,colframe=teal!35!black,title={Before you move on}]
What does \ru{смелость} mean? (\textquotedblleft{}Courage, boldness\textquotedblright{}.) Say it once more.
\end{tcolorbox}
\section[(dialogue)]{Feelings, health and family — a review}

\label{lesson:RU-R476-review-469-476}



\begin{quote}
Say the Russian for jealousy and to be jealous.
\end{quote}

\subsection*{Your turn}
Say these aloud:

\begin{itemize}
\raggedright
  \item jealousy, to be jealous, jealous, envy, to envy
  \item shame, to be ashamed, to get embarrassed, to blush, shy
  \item hurt feelings, to take offence, to hurt, touchy, to forgive
  \item to irritate, irritation, to be angry, fury, quick-tempered
  \item disappointment, disappointed, relief, despair, to sigh
  \item to be nervous, anxious, to worry, worry, to calm down
  \item delight, to admire, to be proud, pride, to boast
  \item a fright, to get frightened, to flinch, a coward, courage
\end{itemize}

\begin{tcolorbox}[breakable,colback=teal!4,colframe=teal!35!black,title={Before you move on}]
Jealousy? (\textbf{\ru{ревность}}.) To be jealous? (\textbf{\ru{ревновать}}.) Jealous? (\textbf{\ru{ревнивый}}.) Envy? (\textbf{\ru{зависть}}.) To envy? (\textbf{\ru{завидовать}}.) Shame? (\textbf{\ru{стыд}}.) To be ashamed? (\textbf{\ru{стыдиться}}.) To get embarrassed? (\textbf{\ru{смущаться}}.) To blush? (\textbf{\ru{краснеть}}.) Shy? (\textbf{\ru{застенчивый}}.) Hurt feelings? (\textbf{\ru{обида}}.) To take offence? (\textbf{\ru{обижаться}}.) To hurt? (\textbf{\ru{обижать}}.) Touchy? (\textbf{\ru{обидчивый}}.) To forgive? (\textbf{\ru{прощать}}.) To irritate? (\textbf{\ru{раздражать}}.) Irritation? (\textbf{\ru{раздражение}}.) To be angry? (\textbf{\ru{злиться}}.) Fury? (\textbf{\ru{ярость}}.) Quick-tempered? (\textbf{\ru{вспыльчивый}}.) Disappointment? (\textbf{\ru{разочарование}}.) Disappointed? (\textbf{\ru{разочарованный}}.) Relief? (\textbf{\ru{облегчение}}.) Despair? (\textbf{\ru{отчаяние}}.) To sigh? (\textbf{\ru{вздыхать}}.) To be nervous? (\textbf{\ru{нервничать}}.) Anxious? (\textbf{\ru{тревожный}}.) To worry? (\textbf{\ru{беспокоиться}}.) Worry? (\textbf{\ru{беспокойство}}.) To calm down? (\textbf{\ru{успокаиваться}}.) Delight? (\textbf{\ru{восторг}}.) To admire? (\textbf{\ru{восхищаться}}.) To be proud? (\textbf{\ru{гордиться}}.) Pride? (\textbf{\ru{гордость}}.) To boast? (\textbf{\ru{хвастаться}}.) A fright? (\textbf{\ru{испуг}}.) To get frightened? (\textbf{\ru{пугаться}}.) To flinch? (\textbf{\ru{вздрагивать}}.) A coward? (\textbf{\ru{трус}}.) Courage? (\textbf{\ru{смелость}}.)
\end{tcolorbox}
